\subsubsection{The lost asset and the obligation still to be met}
\label{sec:onerous-contract-sequence}

Writing down a contract-cost asset does not necessarily recognize the whole
loss on the contract. The asset represents expenditure already incurred.
The company may still have to spend more to perform than the customer will
pay. IAS~37.66--69 requires a provision for that remaining onerous obligation,
after recognizing impairment of assets used in fulfilling the contract
\citep{eu20231803}.

Suppose a company has cash~200, an eligible unamortized contract-cost asset~20
and capital~220. One service remains to be delivered at completion. The customer
will pay~100; incremental future cash costs are~100, and an allocation of other
directly related contract cash costs adds~30. The existing asset's amortization
is separate from those future payments. There are no advances, taxes, foreign
exchange or reimbursements, and the time value of money is immaterial. Assume
the required other-asset impairment tests have been completed and the later
CGU assessment finds no additional loss. These are necessary supplied facts,
not consequences of the calculation.

Future direct costs~130 exceed the remaining consideration~100, so none of the
existing cost asset is recoverable. Recognize impairment~20 first. If the
company fulfills the contract, a further cash loss of \(130-100=30\) remains.
Assume it can instead cancel by paying enforceable compensation~50, with no
remaining benefit or obligation. The least unavoidable loss is then~30, giving
a provision~30. The total immediate loss is \(20+30=50\), and equity is
\(220-50=170\). Cash~200 less provision~30 confirms that amount.

The distinction also gives a compact calculation. Let \(A_0\) be original
unamortized contract cost, \(B\) the remaining collectible consideration,
\(D\) the complete direct future cash cost, and \(K\) the available
cancellation payment. Before any service is delivered, the asset after its
test, the remaining fulfilment loss and the provision are
\begin{align}
 A'&=\min\{A_0,\max(B-D,0)\},\\
 F&=\max\{D+A'-B,0\},\qquad
 P=\min\{F,K\}.
\end{align}
Here \(A'=0\), \(F=30\) and \(P=30\). The asset loss \(A_0-A'\)
has already been recognized and does not enter the provision again. A
provision of~50 obtained by adding the old asset~20 to the remaining cash
loss~30 would charge the asset loss twice. Comparing the gross cost~130 with
the cancellation payment~50 before subtracting revenue~100 would also be
wrong: cancellation does not earn the performance receipt.

The cost population matters as much as the order. IAS~37.68A includes both
incremental costs and allocations of costs directly related to fulfilling
contracts. Ignoring the allocated~30 would incorrectly produce no provision.
That paragraph includes depreciation among possible allocations; this example
uses cash costs to keep the cash reconciliation visible. Unrelated general
administration is excluded from these direct costs. If the company later
pays~15 of such administration, it records a separate period expense and
cannot use this contract's provision to absorb it.

Now complete the service, receive~100 and pay~130. Use provision~30 against
the expenditure for which it was recognized: debit the provision~30 and
current service expense~100, and credit cash~130. Recognize revenue and the
cash receipt~100. There is no additional completion loss. Cash is
\(200+100-130=170\), the cost asset and provision are zero, and equity
remains~170. This use of the provision follows IAS~37.61--62; it is not
permission to release provisions against unrelated expenses. With the separate
administration payment, final cash and equity would instead be~155.

A cheaper exit changes the amount recognized before management acts. If
cancellation costs~20, the provision is~20 and total immediate loss is~40.
Cancelling then pays~20 and leaves cash and equity~180. Management could still
choose to fulfill for commercial reasons. In that history, performance creates
an additional loss~10 and final cash is again~170. Measuring the minimum
unavoidable loss does not choose the company's action. A freely cancellable
contract has no onerous provision, although a worthless cost asset still
requires its own write-down.

Estimates must be kept current. Suppose the original future cost~130 later
falls to~90 before any service occurs. Consideration~100 now supports asset~10:
reverse~10 of its impairment, within the original unamortized ceiling, and
release the provision~30. The cumulative loss becomes~10. Completion then
amortizes asset~10, pays~90 and earns~100, leaving cash and equity~210.
IFRS~15.104 limits the asset reversal, while IAS~37.59 requires the separate
provision review. Partial service, material discounting and further CGU losses
would change this calculation and require their own dated measurements.

